% ==============================================================================
%  PREÂMBULO COMPLEMENTAR (config/preambulo.tex)
% ==============================================================================
%
%  O que este arquivo faz
%    Carrega pacotes e define estilos e comandos usados no texto: tabelas,
%    caixas de destaque, diagramas, listagens de código e cores dos links.
%
%  O que este arquivo NÃO faz
%    A formatação ABNT (margens, fonte, espaçamento, títulos, listas, sumário,
%    referências) está na classe uninove-ppgi.cls. Não a altere aqui.
%
%  Ordem de carregamento (definida em main.tex)
%    1. classe uninove-ppgi.cls
%    2. pacote revisao.sty      (define \ifmodofinal, usado na Seção 6 abaixo)
%    3. este arquivo
%    4. config/dados.tex
%
%  Antes de acrescentar um pacote, verifique se ele já não é carregado pela
%  classe (pacotes repetidos com opções diferentes geram erro "option clash").
%
%  Sumário deste arquivo
%    1. Tabelas e posicionamento de flutuantes
%    2. Atalhos e comandos auxiliares
%    3. Caixas de destaque e de orientação ao aluno
%    4. Diagramas em TikZ
%    5. Listagens de código e expressões de busca
%    6. Cores dos links conforme o modo de compilação
% ==============================================================================


% ==============================================================================
%  1. TABELAS E POSICIONAMENTO DE FLUTUANTES
% ==============================================================================

% float: habilita a opção [H], que fixa figuras, quadros e tabelas
% exatamente no ponto do texto em que foram escritos.
%   Uso: \begin{figure}[H] ... \end{figure}
\usepackage{float}

% longtable: tabelas que continuam na página seguinte.
%   Uso: \begin{longtable}{...} ... \end{longtable}
%   (não deve ser colocada dentro de \begin{table})
\usepackage{longtable}

% Observação: xcolor, colortbl, tabularx, booktabs, multirow e makecell já são
% carregados pela classe. A opção [table] do xcolor é repassada em main.tex.


% ==============================================================================
%  2. ATALHOS E COMANDOS AUXILIARES
% ==============================================================================

% \mes: preenche de cinza uma célula do cronograma (Capítulo 5).
%   Uso: 1 & Atividade & \mes & \mes & & \\ \hline
\newcommand{\mes}{\cellcolor{black!55}}


% ==============================================================================
%  3. CAIXAS DE DESTAQUE E DE ORIENTAÇÃO AO ALUNO
% ==============================================================================

% tcolorbox e mdframed: pacotes para caixas com moldura e fundo colorido.
\usepackage[most]{tcolorbox}
\usepackage{mdframed}

% ------------------------------------------------------------------------------
% Ambiente "destaque": caixa com fundo cinza claro e filete à esquerda.
% Quebra de página automaticamente, se necessário.
%   Uso: \begin{destaque} texto \end{destaque}
% ------------------------------------------------------------------------------
\newmdenv[
  linecolor=gray!50, linewidth=0.4pt,           % cor e espessura do filete
  leftline=true, rightline=false,               % filete apenas à esquerda
  topline=false, bottomline=false,
  backgroundcolor=gray!5,                       % fundo cinza bem claro
  innerleftmargin=10pt, innerrightmargin=6pt,   % espaçamento interno
  innertopmargin=6pt, innerbottommargin=6pt,
  skipabove=\baselineskip, skipbelow=\baselineskip  % espaço antes e depois
]{destaque}

% ------------------------------------------------------------------------------
% Ambiente "orientacao": caixa com instruções ao aluno.
% Aparece no PDF apenas no modo [revisao]; no modo [final], o conteúdo é
% descartado automaticamente, sem necessidade de apagá-lo do código.
%   Uso: \begin{orientacao}[Título opcional]
%          Texto da orientação.
%        \end{orientacao}
% Depende do ambiente "destaque" (acima) e de \ifmodofinal (revisao.sty).
% ------------------------------------------------------------------------------
\ifmodofinal
  % Modo final: tudo entre \begin{orientacao} e \end{orientacao} é ignorado
  \newenvironment{orientacao}[1][]{\comment}{\endcomment}
\else
  % Modo revisão: caixa de destaque com o título "Orientação ao aluno"
  \newenvironment{orientacao}[1][]
  {\begin{destaque}\small\setlength{\parindent}{0pt}%
      \textbf{Orientação ao aluno\if\relax\detokenize{#1}\relax\else: #1\fi.}%
      \par\smallskip}
      {\end{destaque}}
\fi


% ==============================================================================
%  4. DIAGRAMAS EM TIKZ
% ==============================================================================

% Bibliotecas do TikZ (o pacote tikz já é carregado pela classe):
%   shapes.geometric  formas geométricas (losangos, elipses etc.)
%   arrows.meta       pontas de seta configuráveis
%   positioning       posicionamento relativo (below=of, right=of)
%   fit               caixas que envolvem outros nós
%   calc              cálculos com coordenadas
\usetikzlibrary{shapes.geometric, arrows.meta, positioning, fit, calc}

% Estilos padronizados para fluxogramas e desenhos da pesquisa.
%   Uso: \node[box] (a) {Etapa 1};
%        \node[box, below=of a] (b) {Etapa 2};
%        \draw[arrow] (a) -- (b);
\tikzset{
% box: caixa principal, arredondada, sem hifenização no texto interno
box/.style={
    draw, rectangle, rounded corners=4pt,
    minimum width=9cm, minimum height=1.0cm, text width=6.8cm,
    align=center, font=\small, inner sep=6pt, fill=white,
    execute at begin node={\hyphenpenalty=10000\exhyphenpenalty=10000}
  },
% srcbox: caixa compacta, para fontes de dados no topo de um diagrama
srcbox/.style={
    draw, rectangle, rounded corners=4pt,
    minimum width=3.2cm, minimum height=1.4cm, text width=2.9cm,
    align=center, font=\small, inner sep=8pt, fill=white,
    execute at begin node={\hyphenpenalty=10000\exhyphenpenalty=10000}
  },
% arrow: conector com ponta Stealth de 6 pt
arrow/.style={
-{Stealth[length=6pt, width=5pt]},
line width=0.8pt
}
}


% ==============================================================================
%  5. LISTAGENS DE CÓDIGO E EXPRESSÕES DE BUSCA
% ==============================================================================

% upquote: mantém as aspas retas (") nas listagens, em vez de aspas curvas.
% Essencial para que expressões de busca copiadas do PDF funcionem nas bases.
\usepackage{upquote}

% ------------------------------------------------------------------------------
% Estilo "codigo": trechos de programa, com numeração de linhas e cores.
%   Uso: \begin{lstlisting}[style=codigo, language=Python]
%          ...código...
%        \end{lstlisting}
%   Linguagens comuns: Python, R, SQL, Java, C, [LaTeX]TeX.
% ------------------------------------------------------------------------------
\lstdefinestyle{codigo}{
  basicstyle=\footnotesize\ttfamily,           % fonte monoespaçada pequena
  keywordstyle=\color{blue}\bfseries,          % palavras reservadas
  commentstyle=\color{green!50!black},         % comentários
  stringstyle=\color{red},                     % textos entre aspas
  numbers=left, numberstyle=\tiny, numbersep=5pt,  % numeração das linhas
  frame=single, rulecolor=\color{gray},        % moldura
  xleftmargin=20pt, framexleftmargin=15pt,     % recuo para a numeração
  tabsize=3,
  breaklines=true,                             % quebra linhas longas
  showstringspaces=false                       % não marca espaços em textos
}

% ------------------------------------------------------------------------------
% Estilo "busca": expressões de busca em bases de dados (Apêndice A).
% Preserva aspas retas, quebras de linha e recuos, de modo que a expressão
% possa ser copiada do PDF e colada na base sem alterações.
%   Uso: \begin{lstlisting}[style=busca]
%          TITLE-ABS-KEY("termo 1" AND "termo 2")
%        \end{lstlisting}
% ------------------------------------------------------------------------------
\lstdefinestyle{busca}{
basicstyle=\footnotesize\ttfamily,
frame=single, rulecolor=\color{black},
breaklines=true, breakatwhitespace=true,     % quebra apenas em espaços
columns=fullflexible, keepspaces=true,       % espaçamento natural
showstringspaces=false, upquote=true,        % aspas retas
% Acentos dentro das listagens (o listings não os reconhece sozinho)
inputencoding=utf8, extendedchars=true,
literate={á}{{\'a}}1 {é}{{\'e}}1 {í}{{\'i}}1 {ó}{{\'o}}1 {ú}{{\'u}}1
{ã}{{\~a}}1 {õ}{{\~o}}1 {ç}{{\c{c}}}1 {â}{{\^a}}1 {ê}{{\^e}}1
}


% ==============================================================================
%  6. CORES DOS LINKS CONFORME O MODO DE COMPILAÇÃO
% ==============================================================================
% O modo é escolhido em main.tex: \usepackage[revisao]{revisao} ou [final].
%
%   Modo revisão: links em azul (facilita a navegação durante a orientação)
%                 e indicação "Citado na pág." após cada referência, que
%                 revela obras listadas mas não citadas no texto.
%   Modo final:   todo o texto em preto, como exige o manual da biblioteca,
%                 e sem a indicação "Citado na pág.", que não é ABNT.
% ------------------------------------------------------------------------------
\ifmodofinal
  \hypersetup{colorlinks=true, linkcolor=black, citecolor=black,
    filecolor=black, urlcolor=black}
  \renewcommand*{\backrefalt}[4]{}             % remove "Citado na pág."
\else
  \hypersetup{colorlinks=true, linkcolor=blue, citecolor=blue,
    filecolor=blue, urlcolor=blue}
\fi

% ==============================================================================
%  7. METADADOS DO PDF (Arquivo > Propriedades no leitor de PDF)
% ==============================================================================
% Preenchidos automaticamente; não é preciso editar nada aqui.
%   Título, autor e assunto: campos de config/dados.tex
%   Palavras-chave: \PalavrasChave de pre-textuais/resumo.tex
% Como o resumo é lido depois da primeira página, as palavras-chave são
% guardadas no arquivo .aux e gravadas no PDF na compilação seguinte
% (o latexmk e o LaTeX Workshop já compilam mais de uma vez).
% ------------------------------------------------------------------------------
\makeatletter
\providecommand{\pdf@palavraschave}{}
\AtBeginDocument{%
  \hypersetup{
    pdftitle    = {\titulo},
    pdfauthor   = {\autor},
    pdfsubject  = {\tipoexame\ -- \siglaprograma, Universidade Nove de Julho (\ano)},
    pdfkeywords = {\pdf@palavraschave},
    pdfcreator  = {\autor}
  }%
}
\AtEndDocument{%
  \ifdefined\palavraschave
    \immediate\write\@auxout{%
      \gdef\string\pdf@palavraschave{\unexpanded\expandafter{\palavraschave}}}%
  \fi}
\makeatother